% 加藤讲义：原书106--133页；物理111--138页。
\origpage{111}{106}
由 Beilinson（[Be$_2$]）的结果，$\Psi\otimes_A\Sigma^{-1}$ 存在一组 $A$-基，它在映射 $a$ 下于 $A\otimes\mathbb R$ 中的像为
\[
 \pm\lim_{s\to0}s^{-1}L_{A,S}(V^*(1),0),
\]
\footnote{译注：原书此处将极限内 $L$ 函数的第二个变量写为 $0$，原书如此，照录。}
而它在映射 $b$ 下于 $H^1_{\mathrm{lim}}\otimes_A\Sigma^{-1}$ 中的像等于下面这个元素：将 3.3.5 中的 $z_+(X,\mathcal F_+)$ 沿映射
\begin{align*}
&H^1(X,\mathcal F_+)\otimes_\Lambda H^0(\mathbb R,\mathcal F_+(-1))^{-1}\\
&\qquad\longrightarrow H^1(X,\mathcal F)\otimes_\Lambda H^0(\mathbb R,\mathcal F(-1))^{-1}
 \simeq H^1(X,\mathcal F)\otimes_A\Sigma^{-1}
\end{align*}
所得的像。这里 $\Lambda_+$ 和 $\mathcal F_+$ 如 3.3.5 所定义；第一个映射由 $\mathcal F_+\otimes_{\Lambda_+}\Lambda\simeq\mathcal F$ 诱导，其中
\[
 \Lambda_+\longrightarrow\Lambda,\qquad
 \sigma\longmapsto\chi_{\mathrm{cyclo}}(\sigma)\chi_{\mathcal F}(\sigma)^{-1}
 \quad(\sigma\in\Gal(\mathbb Q^{\mathrm{ab},S}/\mathbb Q)).
\]
这就说明，当 $r\ge2$ 且 $\chi_{\mathcal F}(\tau)=-1$ 时，§3 的构造与 3.2.2(v) 相容。当 $r=1$ 且 $\chi_{\mathcal F}(\tau)=-1$ 时，借助 3.3.6 再作更细致的分析，也能得到类似结论，此处不再讨论。

\textbf{注 3.4.10。} 正如在 3.4.6--3.4.8 中所见，§1.3 讨论的 Riemann $\zeta$ 函数的第二个侧面，即其 $p$ 进性质，已包含在我们的\term{岩泽主猜想}{Iwasawa main conjecture}{conjecture principale d'Iwasawa}中。

粗略地说，我们的方法通过如下方式涵盖第二个侧面。若整数 $r,r'$ 满足
\[
 r\equiv r'\pmod{(p-1)p^{n-1}},
\]
则在 $\Spec(\mathbb Z[1/p])$ 上有层的同构 $\mathbb Z/p^n(r)\simeq\mathbb Z/p^n(r')$。因此，$\mathbb Z_p(r)$ 与 $\mathbb Z_p(r')$ 的\term{泽塔元素}{zeta element}{élément zêta}模 $p^n$ 相同，由此产生泽塔值之间的同余关系（1.3.4(3)）。对应 $r\mapsto\mathbb Z_p(r)$ 从 $r\in\mathbb Z$ 延拓到
\[
 r\in\varprojlim_n\mathbb Z/(p-1)p^n\mathbb Z.
\]
这就解释了 Riemann $\zeta$ 函数特殊值的 $p$ 进插值为何存在（1.3.4(4)）。

\origpage{112}{107}
在第三章 §3，我们将讨论岩泽主猜想与 $p$ 进 $L$ 函数之间的一般关系。

\textbf{3.4.11。} 如何构造 $p$ 进泽塔元素显然是一个重要问题，但我还没有什么好主意。能否完全不考虑泽塔值，而只考虑 $p$ 进平展上同调来构造它们呢？

关于构造，我想说的一点是：对于一个 $A$-\term{动机}{motif}{motif} $M$，它的 $p$ 进实现 $M_p$ 的泽塔元素的构造，似乎可以归结到严格 $K$ 理论的情形。思路如下。

设 $S$ 是包含 $p$ 的有限素数集，$A$ 是 $\mathbb Q$ 的若干有限扩张的有限直积，$M$ 是一个 $A$-动机，且 $M_p$ 在 $S$ 之外非分歧。我们希望定义 $z_\Lambda(X,M_p)$，其中
\[
 \Lambda=A\otimes\mathbb Q_p,\qquad X=\Spec(\mathbb Z[1/S]).
\]
预期当 $r\ge0$ 充分大时，对于任意 $n\ge0$，动机
$M(r)\otimes_{\mathbb Q}H^0(\Spec(\mathbb Q(\beta_n)))$ 都是严格 $K$ 理论的；这里 $\beta_n$ 是一个本原 $p^n$ 次单位根。令 $\mathcal O_\Lambda\subset\Lambda$ 如 3.1.2 所定义，并取 $X$ 上的光滑 $\mathcal O_\Lambda$-层 $\mathcal F$，使 $\mathcal F\otimes_{\mathcal O_\Lambda}\Lambda=M_p$。令
\[
 \begin{gathered}
 G=\Gal\bigl((\bigcup_n\mathbb Q(\beta_n))/\mathbb Q\bigr),\qquad
 \Lambda'=\mathcal O_\Lambda[[G]],\\
 \mathcal F'=\varprojlim_n f_{n*}f_n^*(\mathcal F(r))\in D_{\mathrm{ctf}}(X,\Lambda').
\end{gathered}
\]
其中 $f_n$ 是典范态射
$\Spec(\mathbb Z[\beta_n,1/S])\longrightarrow\Spec(\mathbb Z[1/S])$。于是，关于 $\mathcal O_\Lambda$-同态
\[
 \Lambda'\longrightarrow\Lambda,\qquad
 \sigma\longmapsto\chi_{\mathrm{cyclo}}(\sigma)^r\quad(\sigma\in G),
\]
有
\begin{equation*}
 M_p\simeq\mathcal F'\otimes_{\Lambda'}\Lambda.\tag{1}
\end{equation*}
假设对于每个 $n\ge0$，我们已经定义了
$\Lambda[\Gal(\mathbb Q(\beta_n)/\mathbb Q)]$-动机
$M(r)\otimes_{\mathbb Q}H^0(\Spec(\mathbb Q(\beta_n)))$ 的 $p$ 进平展实现
$f_{n*}f_n^*(\mathcal F(r))\otimes_{\mathcal O_\Lambda}\Lambda$ 的泽塔元素。那么，$z_{\Lambda'}(X,\mathcal F')$ 应由这些元素取逆极限而得。
\origpage{113}{108}
而 $z_\Lambda(X,\mathcal F)$ 应是 $z_{\Lambda'}(X,\mathcal F')$ 在 (1) 下的像。

这里有两个问题：
\begin{enumerate}[label=(\roman*)]
\item 在严格 $K$ 理论的情形如何构造泽塔元素？
\item 按上述方法得到的泽塔元素 $z_\Lambda(X,\mathcal F)$ 是否满足 3.2.2(v)？
\end{enumerate}
问题 (ii)，以及与之类似的 3.4.1(1)，可以形象地说成：“卓越的 $p$ 进元素经过 Tate 扭转后，仍然卓越吗？”（“卓越”在这里表示它们与泽塔值有关。）如果我们的岩泽主猜想成立，问题 (ii) 的答案就是肯定的；它是这一主猜想引出的最深刻的问题之一。

问题 (i) 显然同样重要。我唯一能提出的建议，是向 S.~Bloch、A.~Beilinson 以及其他发现了 $K$ 理论中与泽塔值相关的重要元素的人学习，看看怎样构造 $\omega_\gamma\in\det_A(\Psi)$（之后 $z_\gamma$ 就是 $\omega_\gamma$ 在 Chern 类映射下的像）；参见 [Bl] 第 8、9 讲、[Be$_1$] 等。在我看来，发现这类重要元素的唯一已知一般方法，就是张开嘴，等着它们从天上掉下来。不过，我不明白，那些嘴巴小小的人怎么能屡屡接到这样的元素。

\textbf{3.4.12。} $\Delta_\Lambda(X,\mathcal F)$ 的 $\Lambda$-基有很多，但我们的猜想断言，其中存在一个特别重要的基。

这样的特殊基应当有一个优美的定义，正如 §3.3 的情形，也如第三章 §1.2 中某些泽塔元素来自定义优美的\term{椭圆单位}{elliptic unit}{unité elliptique}。已知的 $K$ 理论中与泽塔值有关的重要元素，其定义也都很优美（参见 [Ra]）。该如何理解这一点呢？
\origpage{114}{109}
这种特殊元素的存在，仿佛对一个男子（或女子）而言，总有一位特别的女子（或男子）存在。这样的类比，能帮助我们寻找泽塔元素吗？

\textbf{3.4.13。} 我想说明，泽塔元素一旦被发现，其力量很大，对\term{算术几何}{arithmetic geometry}{géométrie arithmétique}的研究十分有用。

例如，设 $X$ 是 $\mathbb Q$ 上一个 $g$ 维\term{阿贝尔簇}{abelian variety}{variété abélienne}，$S$ 是包含 $p$ 的有限素数集，且 $X$ 在 $S$ 之外有好约化。假设 $L_S(H^1(X),s)$ 能亚纯延拓到整个 $\mathbb C$，并且在 $s=1$ 处的阶为 $0$。令
\[
 V=H^1(X\otimes_{\mathbb Q}\overline{\mathbb Q},\mathbb Q_p)(1).
\]
那么 §2.3 的猜想，或三元组 $(\Spec(\mathbb Z[1/S]),\mathbb Q_p,V)$ 的泽塔元素所应满足的性质 3.2.2(v)，断言存在
\[
 z\in\bigwedge^g H^1(\mathbb Z[1/S],V),
\]
使其在 $\bigwedge^g\exp^*$ 下，于一维 $\mathbb Q_p$-向量空间
$\bigwedge^g\Gamma(X,\Omega^1_{X/\mathbb Q})\otimes\mathbb Q_p$ 中的像非零。假设我们已经找到了这样的 $z$。（事实上，在一些十分特殊的情形中，我们将在即将发表的论文 [Ka$_3$] 中，利用 3.4.11 的想法构造 $z$。）那么如下所示，我们能推出一个很强的结论：$X(\mathbb Q)$ 是有限集。因此，我认为，泽塔元素这一思想对于算术几何的研究，是一个有用的工作假设。

现在从 $z$ 的存在推出 $X(\mathbb Q)$ 的有限性。映射
\[
 H^1(\mathbb Z[1/S],V)\longrightarrow\Gamma(X,\Omega^1_{X/\mathbb Q})\otimes\mathbb Q_p
\]
是满射。考虑交换图
\[
\begin{tikzcd}[ampersand replacement=\&,column sep=small]
X(\mathbb Q)\otimes\mathbb Q\arrow[r,hook]\arrow[d]\&X(\mathbb Q_p)\otimes\mathbb Q\arrow[d]\&\operatorname{Lie}(X)\otimes\mathbb Q_p\arrow[l,"{\exp}"',"{\simeq}"]\arrow[d]\\
H^1(\mathbb Q,V_pX)\arrow[r]\&H^1(\mathbb Q_p,V_pX)\arrow[r,"{(*)}"]\&\Hom_{\mathbb Q_p}(H^1(\mathbb Z[1/S],V),\mathbb Q_p).
\end{tikzcd}
\]
左边和中间的竖箭头按 2.3.6 定义。
\origpage{115}{110}
箭头 $(*)$ 则由\term{Tate 对偶}{Tate duality}{dualité de Tate}定义。右边的竖箭头是上述满射的 $\mathbb Q_p$-对偶，因而是单射。另一方面，由 Tate 对偶的整体互反律，下行两映射的复合为零。因此
$X(\mathbb Q)\otimes\mathbb Q\longrightarrow X(\mathbb Q_p)\otimes\mathbb Q$ 为零，从而 $X(\mathbb Q)$ 有限。

\textbf{3.4.14。} 为使叙述简洁，本节一直排除了 $p=2$。对于 $p=2$，我们的岩泽主猜想可以如下表述。考虑 3.2.1 中的三元组 $(X,\Lambda,\mathcal F)$，并只考虑 $X=\Spec(\mathbb Z[1/p])$ 的情形。（回忆在 3.2.2 中，利用 3.2.2(ii)，一般情形的泽塔元素可归结到这一情形。）我们如下定义 $D_{\mathrm{ctf}}(\Lambda)$ 中的对象 $R\Gamma_c(X,\mathcal F)$。先假设 $\Lambda$ 是有限环来叙述定义；一般情形显然可归结于此。取 $\mathcal F$ 的\term{内射消解}{injective resolution}{résolution injective} $\mathcal F\longrightarrow I^\bullet$，令 $R\Gamma_c(X,\mathcal F)$ 为
\begin{equation*}
 \Gamma(X,I^\bullet)\longrightarrow\Gamma(\mathbb R,I^\bullet)\oplus\Gamma(\mathbb Q_p,I^\bullet)\tag{1}
\end{equation*}
的\term{映射纤维}{mapping fiber}{fibre d'un morphisme}，即双重复形 (1) 所伴随的单复形。定义
\[
 \Delta_{c,\Lambda}(X,\mathcal F)=\{\det_\Lambda R\Gamma_c(X,\mathcal F)\}^{-1}.
\]
令 $\mathcal F^*=R\Hom_\Lambda(\mathcal F,\Lambda)$。当 $p\ne2$ 时，代数数论中平展上同调的对偶理论（[Ma]）给出典范的 $\Lambda$-同构
\[
 \Delta_{c,\Lambda}(X,\mathcal F)\simeq\Delta_\Lambda(X,\mathcal F^*(1)).
\]
因此在 $p\ne2$ 时，可以利用 $\Delta_{c,\Lambda}(X,\mathcal F^*(1))$ 改写三元组 $(X,\Lambda,\mathcal F)$ 的 3.2.2。再利用 $\Delta_{c,\Lambda}(X,\mathcal F)$，就可自然地写出包括 $p=2$ 在内的 3.2.2 的推广。

在 §3 正文中，我更愿意用 $\Delta$ 而非 $\Delta_c$，因为如前所见，有些卓越的元素栖身于 $R\Gamma$ 的 $H^1$ 中，却不在 $R\Gamma_c$ 的上同调中。

\origpage{116}{111}
\section*{§3.5. 有限域上的簇}
\addcontentsline{toc}{section}{§3.5. 有限域上的簇}
\textbf{3.5.1。} 回忆经典岩泽理论的出发点，是岩泽对有限域上曲线与数域整数环的谱之间深刻类比的如下信念：前者的泽塔函数可用 Jacobian 簇的 Tate 模，或用平展上同调理论表示，后者的泽塔函数也应有相应的表示。

在 §3.5 中，令 $(X,\Lambda,\mathcal F)$ 为 3.2.2(iv) 中的三元组，即几何情形，$\Lambda$ 如 3.1.1(2) 所定义。本节的目的是说明，在 3.5.3 给出的温和条件下，3.2.2(iv) 定义的泽塔元素 $z_\Lambda(X,\mathcal F)$ 与泽塔值之间的关系确实成立，而不仅是猜想（命题 3.5.7）。例如，若所有 $m\in\mathbb Z$ 都有 $H^m(X\otimes_{\mathbb F_\ell}\overline{\mathbb F}_\ell,\mathcal F)=0$，则命题 3.5.7 说明，典范等同
\[
 \Delta_\Lambda(X,\mathcal F)=\{\det_\Lambda(\{0\})\}^{-1}=\Lambda
\]
将 $z_\Lambda(X,\mathcal F)$ 送到泽塔函数 $Z_\Lambda(X,D_\Lambda(\mathcal F),t)$ 在 $t=1$ 处的值，至多相差一个符号。这里
\[
 D_\Lambda(\mathcal F)=R\Hom_\Lambda(\mathcal F,f^!\Lambda),
\]
$f$ 是结构态射 $X\longrightarrow\Spec(\mathbb F_\ell)$，泽塔函数定义为
\begin{align*}
 Z_\Lambda(X,\mathcal F,t)
 &=\prod_{q,x}\det_\Lambda\bigl(1-\sigma_x^{-1}t^{\deg(x)};\mathcal H^q(\mathcal F)_{\bar x}\bigr)^{(-1)^{q-1}}\in\Lambda[[t]]\\
 &=\det_\Lambda\bigl(1-\sigma_\ell^{-1}t;R\Gamma_c(X\otimes_{\mathbb F_\ell}\overline{\mathbb F}_\ell,\mathcal F)\bigr)^{-1}.
\end{align*}
其中 $q$ 遍历整数，$x$ 遍历 $X$ 的闭点，$\mathcal H^q(\mathcal F)$ 是 $\mathcal F$ 的第 $q$ 个上同调层，$(\ )_{\bar x}$ 表示 $x$ 上方几何点处的茎，$\sigma_\ell\in\Gal(\overline{\mathbb F}_\ell/\mathbb F_\ell)$ 是\term{算术 Frobenius}{arithmetic Frobenius}{Frobenius arithmétique}；最后一个等式来自 Grothendieck [Gr]。

\origpage{117}{112}
\textbf{3.5.2。} 令 $\mathcal C(\mathcal F)$ 为复形
\[
 H^0(X,\mathcal F)\xrightarrow{\cup\theta}H^1(X,\mathcal F)\xrightarrow{\cup\theta}H^2(X,\mathcal F)\xrightarrow{\cup\theta}\cdots,
\]
其中 $\cup\theta$ 表示与
\[
 \theta\in H^1(\mathbb F_\ell,\mathbb Z_p)=\Hom_{\mathrm{cont}}(\Gal(\overline{\mathbb F}_\ell/\mathbb F_\ell),\mathbb Z_p)
\]
的杯积；此生成元将 $\sigma_\ell$ 送到 $1\in\mathbb Z_p$。对 $i\in\mathbb Z$，令 $P^i$、$Q^i$ 分别为
\[
 1-\sigma_\ell:H^i(\overline X,\mathcal F)\longrightarrow H^i(\overline X,\mathcal F)
\]
的核与余核。在\term{谱序列}{spectral sequence}{suite spectrale}
\[
 E_2^{i,j}=H^i(\mathbb F_\ell,H^j(\overline X,\mathcal F))\Longrightarrow E_\infty^{i+j}=H^{i+j}(X,\mathcal F)
\]
中，$E_2^{i,j}$ 在 $i=0$ 时为 $P^j$，在 $i=1$ 时为 $Q^j$，在 $i\ge2$ 时为零。因此有短正合列
\begin{equation*}
 0\longrightarrow Q^{i-1}\longrightarrow H^i(X,\mathcal F)\longrightarrow P^i\longrightarrow0
 \qquad(i\in\mathbb Z).\tag{1}
\end{equation*}
而 $\cup\theta$ 等于复合映射
\[
 H^i(X,\mathcal F)\longrightarrow P^i\xrightarrow{\alpha_i}Q^i\longrightarrow H^{i+1}(X,\mathcal F),
\]
其中首尾两箭头来自 (1)，$\alpha_i$ 则是
\[
 P^i\xrightarrow{\mathrm{incl}}H^i(\overline X,\mathcal F)\xrightarrow{\mathrm{proj}}Q^i
\]
的复合，$\mathrm{incl}$ 为包含映射，$\mathrm{proj}$ 为商映射。因此有：
\begin{statement}{引理 3.5.3}
下列两个条件等价：
\begin{enumerate}[label=(\arabic*)]
\item 复形 $\mathcal C(\mathcal F)$ 正合；
\item 对所有 $i$，映射 $\alpha_i:P^i\longrightarrow Q^i$ 均为双射。
\end{enumerate}
\end{statement}
\textbf{注 3.5.4。} 若所有 $i$ 上 Frobenius $\sigma_\ell$ 对 $H^i(\overline X,\mathcal F)$ 的作用都是半单的，则上述条件 (2)，从而条件 (1)，均成立。预期当 $\mathcal F$ 来自“$X$ 上的一个动机”时，这一半单性成立。

以下引理由 $\mathcal F$ 的 $R\Gamma$ 与 $D_X(\mathcal F)$ 的 $R\Gamma_c$ 之间的 Poincaré 对偶得到。
\origpage{118}{113}
\begin{statement}{引理 3.5.5}
假设 3.5.3 的等价条件成立。则 $Z_\Lambda(X,D_\Lambda(\mathcal F),t)$ 在 $t=1$ 处的阶 $e$，作为 $\Spec(\Lambda)\longrightarrow\mathbb Z$ 的函数，等于\translatornote{10}
\[
 \sum_{i\in\mathbb Z}(-1)^i\rank_\Lambda(P^i)=\sum_{i\in\mathbb Z}(-1)^i\rank_\Lambda(Q^i).
\]
\end{statement}
\textbf{3.5.6。} 假设 3.5.3 的等价条件成立，定义同构
\[
 \lambda:\{\det_\Lambda(R\Gamma(X,\mathcal F))\}^{-1}\longrightarrow\Lambda
\]
为下列映射的复合：
\begin{align*}
 \{\det_\Lambda(R\Gamma(X,\mathcal F))\}^{-1}
 &\simeq\bigotimes_{i\in\mathbb Z}\{\det_\Lambda(H^i(X,\mathcal F))\}^{(-1)^{i-1}}\\
 &\simeq\{\det_\Lambda(\mathcal C(\mathcal F))\}^{-1}\simeq\Lambda.
\end{align*}
最后一个同构来自 $\mathcal C(\mathcal F)$ 的无环性，即它在导出范畴中等于零。一般说来，这一同构 $\lambda$ 与 3.2.2(vi) 中讨论的同构不同；二者之差正是命题 3.5.7 的主题。若 $R\Gamma(X,\mathcal F)=0$，那么由 $\det_\Lambda(0)\simeq\Lambda$ 得到的那个最简单的同构就是 $\lambda$。

除去符号，$\lambda$ 等于下列同构的复合：
\begin{align*}
 \{\det_\Lambda(R\Gamma(X,\mathcal F))\}^{-1}
 &\simeq\bigotimes_{i\in\mathbb Z}\{\det_\Lambda(H^i(X,\mathcal F))\}^{(-1)^{i-1}}\\
 &\simeq\bigotimes_{i\in\mathbb Z}\{\det_\Lambda(P^i)\otimes(\det_\Lambda(Q^i))^{-1}\}^{(-1)^{i-1}}
 \simeq\Lambda.
\end{align*}
第二个同构由短正合列 3.5.2(1) 诱导，最后一个同构由 $\alpha_i:P^i\xrightarrow{\sim}Q^i$ 诱导。
\begin{statement}{命题 3.5.7}
假设 3.5.3 的等价条件成立。则同构 $\lambda$ 将 $z_\Lambda(X,\mathcal F)$ 送到
\[
 \pm\left.\{(1-t)^{-e}Z_\Lambda(X,D_\Lambda(\mathcal F),t)\}\right|_{t=1}\in\Lambda^\times,
\]
其中 $e:\Spec(\Lambda)\longrightarrow\mathbb Z$ 是 $Z_\Lambda(X,D_\Lambda(\mathcal F),t)$ 在 $t=1$ 处的阶。
\end{statement}

\origpage{119}{114}
该命题由以下引理得到。
\begin{statement}{引理 3.5.8}
设 $R$ 是有限多个域的直积，$C$ 是 $R$-模的\term{完美复形}{perfect complex}{complexe parfait}，$h:C\longrightarrow C$ 为同态。对 $i\in\mathbb Z$，令 $P^i$、$Q^i$ 分别为
$H^i(1-h):H^i(C)\longrightarrow H^i(C)$ 的核与余核，并假设对每个 $i$，复合映射
\[
 \alpha_i:P^i\xrightarrow{\mathrm{incl}}H^i(C)\xrightarrow{\mathrm{proj}}Q^i
\]
为双射。令 $e:\Spec(R)\longrightarrow\mathbb Z$ 为
\[
 f(t):=\prod_{q\in\mathbb Z}\det_R(1-ht;H^q(C)\longrightarrow H^q(C))^{(-1)^q}
\]
在 $t=1$ 处的阶。则
\[
 e=\sum_{i\in\mathbb Z}(-1)^i\rank_R(P^i)=\sum_{i\in\mathbb Z}(-1)^i\rank_R(Q^i),
\]
而复合映射
\begin{align*}
 R&\simeq\det_R(C)\otimes_R\det_R(C)^{-1}\\
 &\simeq\bigotimes_{i\in\mathbb Z}\bigl\{\det_R(H^i(C))^{(-1)^i}\otimes_R\det_R(H^i(C))^{(-1)^{i-1}}\bigr\}\\
 &\simeq\bigotimes_{i\in\mathbb Z}\bigl\{\det_R(P^i)^{(-1)^i}\otimes_R\det_R(Q^i)^{(-1)^{i-1}}\bigr\}\simeq R\tag{*}
\end{align*}
将 $1$ 送到 $\pm\lim_{t\to0}(1-t)^{-e}f(t)$。\footnote{译注：原书此处写 $t\to0$，而本引理前文讨论的是 $t=1$ 处的阶；原书如此，照录。}这里在 $(*)$ 中，第一个同构来自两个 $C$ 之间的恒等映射，第二个来自
$\det_R(C)\simeq\bigotimes_{i\in\mathbb Z}\det_R(H^i(C))^{(-1)^i}$，第三个来自正合列
\[
 0\longrightarrow P^i\longrightarrow H^i(C)\xrightarrow{1-h}H^i(C)\longrightarrow Q^i\longrightarrow0,
\]
最后一个则来自同构 $\alpha_i$。
\end{statement}
\begin{proof}
可以归结到 $C$ 集中于零次的情形，并进一步归结到 $1-h$ 是同构或零映射的情形。
\end{proof}
\textbf{注 3.5.9。} 考虑三元组 $(X,\mathbb Z_p,\mathcal F)$，其中 $X$ 如 3.2.2(iv) 所定义。假设 $\mathcal F\otimes_{\mathbb Z_p}\mathbb Q_p$ 满足 3.5.3 的等价条件。那么将 3.5.7 用于 $\mathcal F\otimes_{\mathbb Z_p}\mathbb Q_p$，可知
\[
 \left.\{(1-t)^{-e}Z_{\mathbb Q_p}(X,D_\Lambda(\mathcal F\otimes_{\mathbb Z_p}\mathbb Q_p),t)\}\right|_{t=1}
\]
等于
$\prod_{q\in\mathbb Z}\#(H^q(\mathcal C(\mathcal F)))^{(-1)^q}$ 乘上一个 $\mathbb Z_p^\times$ 中的元素。

\origpage{120}{115}
\section*{第二章\quad 局部理论}
\addcontentsline{toc}{section}{第二章\quad 局部理论}
回忆我们的基本思想：在 Hasse--Weil $L$ 函数的岩泽理论中，与 $B_{\mathrm{dR}}$ 有关的 $p$ 进周期理论，是最重要的局部理论。

本章中，$K$ 表示一个特征为零的\term{完备离散赋值域}{complete discrete valuation field}{corps de valuation discrète complet}，其剩余域 $k$ 是特征 $p>0$ 的完美域。记 $\mathcal O_K$ 为 $K$ 的赋值环，$\mathfrak m_K$ 为其极大理想。用 $K_0$ 表示 $p$-Witt 向量环 $W(k)$ 的分式域，并以典范方式将 $K_0$ 看作 $K$ 的子域。固定 $K$ 的一个代数闭包 $\overline K$，用 $K_{\mathrm{ur}}$ 表示 $\overline K$ 中 $K$ 的最大非分歧扩张。$H^q(K,\ )$ 表示连续 Galois 上同调
$H^q_{\mathrm{cont}}(\Gal(\overline K/K),\ )$（[Ta$_2$]）。令 $\mathbb Z_p(1)=H^0(\overline K,\mathbb Z_p(1))$。

\subsection*{§1. $p$ 进域的 $p$ 进 Galois 表示}
\addcontentsline{toc}{subsection}{§1. $p$ 进域的 $p$ 进 Galois 表示}
本节回顾与 $B_{\mathrm{dR}}$ 有关的 $p$ 进周期理论，不过 §1.2 和 §1.4 包含一个新内容，即“对偶指数映射”。参见 Fontaine [Fo$_1$]、Fontaine--Messing [FM]、Faltings [Fa]、Bloch--Kato [BK]、Illusie [Il]。

\subsection*{§1.1. $B_{\mathrm{dR}}$ 与 de Rham 表示}
我们回顾域 $B_{\mathrm{dR}}$ 和 Fontaine 的 de Rham 表示理论。用\term{晶体上同调}{crystalline cohomology}{cohomologie cristalline}来定义 $B_{\mathrm{dR}}$ 的想法，来自 [FM]，与 [Fo$_1$] 的原始方法不同；这一说明也适用于 §1.3 中 $B_{\mathrm{crys}}$ 的定义。
\origpage{121}{116}
\textbf{1.1.1. $B_{\mathrm{dR}}$ 的定义。} 对 $n\ge1$，令
\[
 B_{n,\overline K/K}=H^0(\Spec(\mathcal O_{\overline K}/p^n)/\Spec(\mathcal O_K/p^n)_{\mathrm{crys}},\mathcal O_{\mathrm{crys}}),
\]
其中 $\Spec(\mathcal O_{\overline K}/p^n)/\Spec(\mathcal O_K/p^n)_{\mathrm{crys}}$ 是概形 $\Spec(\mathcal O_{\overline K}/p^n)$ 在底概形 $\Spec(\mathcal O_K/p^n)$ 上，关于 $p(\mathcal O_K/p^n)$ 上典范\term{分幂结构}{divided power structure}{structure à puissances divisées}的晶体景，$\mathcal O_{\mathrm{crys}}$ 是此景的结构层。于是典范映射
$B_{n,\overline K/K}\longrightarrow\mathcal O_{\overline K}/p^n$ 为满射。令 $J_{n,\overline K/K}$ 为其核，并对 $r\ge0$，将 $B_{n,\overline K/K}$ 的理想 $J^{[r]}_{n,\overline K/K}$ 定义为 $J_{n,\overline K/K}$ 的第 $r$ 个分幂。于是对 $r\ge0$，
\[
 J^{[r]}_{n,\overline K/K}/J^{[r+1]}_{n,\overline K/K}
\]
是秩一的自由 $\mathcal O_{\overline K}/p^n$-模；对 $\mathcal O_K$ 的任一素元 $\pi$，以及 $B_{n,\overline K/K}$ 中任一满足 $x^{p^n}$ 在 $\mathcal O_{\overline K}/p^n$ 中的像等于 $\pi$ 的类的 $x$，它由 $(x^{p^n}-\pi)^{[r]}$ 的类生成。令
\[
 B_{\infty,\overline K/K}=\varprojlim_n B_{n,\overline K/K},\qquad
 J^{[r]}_{\infty,\overline K/K}=\varprojlim_n J^{[r]}_{n,\overline K/K}.
\]
定义
\[
 B^+_{\mathrm{dR}}=\varprojlim_r\bigl((B_{\infty,\overline K/K}/J^{[r]}_{\infty,\overline K/K})\otimes\mathbb Q\bigr).
\]
则 $B^+_{\mathrm{dR}}$ 是完备离散赋值环，而且对 $r\ge0$，典范满射
\[
 B^+_{\mathrm{dR}}\longrightarrow(B_\infty/J_\infty^{[r]})\otimes\mathbb Q
\]
的核等于 $B^+_{\mathrm{dR}}$ 的极大理想的 $r$ 次幂。

将 $B_{\mathrm{dR}}$ 定义为 $B^+_{\mathrm{dR}}$ 的分式域。它因而是完备离散赋值域，$\Gal(\overline K/K)$ 在其上的作用保持赋值。注意 $B^+_{\mathrm{dR}}$ 有一个 $K$-代数结构。已知对于 $K$ 的任意有限扩张 $L$，从 $K$ 的 $B_{\mathrm{dR}}$ 到 $L$ 的 $B_{\mathrm{dR}}$ 的典范映射都是同构。我们将二者等同。
\origpage{122}{117}
\textbf{1.1.2. 函子 $D_{\mathrm{dR}}$。} $D_{\mathrm{dR}}$ 从带有 $\Gal(\overline K/K)$ 连续作用的有限维 $\mathbb Q_p$-向量空间范畴，映到带有以 $\mathbb Z$ 为指标的降滤过的有限维 $K$-向量空间范畴。它定义为
\[
 D_{\mathrm{dR}}(V)=H^0(K,V\otimes_{\mathbb Q_p}B_{\mathrm{dR}}),
\]
其中 $\Gal(\overline K/K)$ 在张量积上作对角作用；其滤过为
\[
 D^i_{\mathrm{dR}}(V)=H^0(K,V\otimes_{\mathbb Q_p}B^i_{\mathrm{dR}})\qquad(i\in\mathbb Z),
\]
$B^i_{\mathrm{dR}}$ 表示 $B_{\mathrm{dR}}$ 中归一化加法赋值至少为 $i$ 的元素组成的部分。

\textbf{1.1.3. de Rham 表示。} 设 $V$ 是带有 $\Gal(\overline K/K)$ 连续作用的有限维 $\mathbb Q_p$-向量空间。已知
\[
 \dim_KD_{\mathrm{dR}}(V)\le\dim_{\mathbb Q_p}V.
\]
若等号成立，则称 $V$ 为 $\Gal(\overline K/K)$ 的一个\term{de Rham 表示}{de Rham representation}{représentation de de Rham}。此时典范映射
\[
 D_{\mathrm{dR}}(V)\otimes_K B_{\mathrm{dR}}\longrightarrow V\otimes_{\mathbb Q_p}B_{\mathrm{dR}}
\]
为双射，并保持滤过；左端的 $\mathrm{fil}^i$ 为 $\sum_{r+s=i}\mathrm{fil}^r\otimes\mathrm{fil}^s$，右端的 $\mathrm{fil}^i$ 为 $V\otimes\mathrm{fil}^i$。de Rham 表示对张量积、外幂、对偶和直和运算稳定，且 $D_{\mathrm{dR}}$ 在 de Rham 表示上与这些运算可交换。de Rham 表示的子表示和商表示仍是 de Rham 的。
\origpage{123}{118}
若 $V$ 是 de Rham 表示，$L$ 是 $\overline K$ 中 $K$ 的有限扩张，则将 $V$ 看作 $\Gal(\overline K/L)$ 的表示时，它仍是 de Rham 的，且
\[
 D_{\mathrm{dR}}(V)\otimes_K L\longrightarrow D_{\mathrm{dR}}(L,V),\qquad
 D^i_{\mathrm{dR}}(V)\otimes_K L\longrightarrow D^i_{\mathrm{dR}}(L,V)
\]
均为同构；这里 $D_{\mathrm{dR}}(L,\ )$ 表示用于 $\Gal(\overline K/L)$ 的表示的 $D_{\mathrm{dR}}$ 函子。

以下 1.1.4 由 Fontaine [Fo$_1$] 提出猜想，Fontaine [Fo$_1$] 和 Fontaine--Messing [FM] 曾在一些特殊情形下证明它。
\begin{statement}{定理 1.1.4（Faltings [Fa]）}
设 $X$ 是 $K$ 上的光滑簇，$m\in\mathbb Z$。则 $p$ 进平展上同调 $H^m_{\et}(X\otimes_K\overline K,\mathbb Q_p)$ 是 $\Gal(\overline K/K)$ 的 de Rham 表示，而且作为滤过 $K$-向量空间，
\[
 D_{\mathrm{dR}}(H^m(X\otimes_K\overline K,\mathbb Q_p))
\]
典范同构于\term{de Rham 上同调}{de Rham cohomology}{cohomologie de de Rham}
$H^m_{\mathrm{dR}}(X/K)=H^m(X,\Omega^\bullet_{X/K})$。
\end{statement}
这里 $H^m_{\mathrm{dR}}(X/K)$ 的滤过按 [De$_1$] 定义。（若 $X$ 在 $K$ 上固有，则
\[
 \mathrm{fil}^r H^m_{\mathrm{dR}}(X/K)=\im\bigl(H^m(X,\Omega^{\ge r}_{X/K})\longrightarrow H^m_{\mathrm{dR}}(X/K)\bigr).
\]
）

\textbf{1.1.5。} 最后回顾 $D_{\mathrm{dR}}$ 与\term{Tate 扭转}{Tate twist}{torsion de Tate}之间的关系。设有限维 $\mathbb Q_p$-向量空间 $V$ 带有 $\Gal(\overline K/K)$ 的连续作用，$r\in\mathbb Z$。则作为 $K$-向量空间，$D_{\mathrm{dR}}(V(r))$ 与 $D_{\mathrm{dR}}(V)$ 等同，但滤过为
\[
 D^i_{\mathrm{dR}}(V(r))=D^{i+r}_{\mathrm{dR}}(V).
\]
此等同如下给出。存在一个典范的 $\mathbb Q_p[\Gal(\overline K/K)]$-同态
$\varepsilon:\mathbb Q_p(1)\longrightarrow B_{\mathrm{dR}}$（见下文），使得 $\mathbb Q_p(1)$ 中任一非零元素的像都是 $B^+_{\mathrm{dR}}$ 的素元。双射
\[
 \mathbb Q_p(r)\otimes_{\mathbb Q_p}B_{\mathrm{dR}}\xrightarrow{\sim}B_{\mathrm{dR}},\qquad
 x^{\otimes r}\otimes y\longmapsto\varepsilon(x)^r y
 \quad(x\in\mathbb Q_p(1),\ x\ne0,\ y\in B_{\mathrm{dR}})
\]
诱导双射 $D_{\mathrm{dR}}(V(r))\longrightarrow D_{\mathrm{dR}}(V)$，我们就以此为等同。

下面回顾 $\varepsilon$ 的定义。
\origpage{124}{119}
此映射由以下同态 $\varepsilon:\mathbb Z_p(1)\longrightarrow B_{\infty,\overline K/K}$ 诱导。令 $\beta=(\beta_n)_{n\ge1}\in\mathbb Z_p(1)$，其中 $\beta_n$ 是 $\overline K$ 中的 $p^n$ 次单位根，满足 $(\beta_{n+1})^p=\beta_n$。对每个 $n\ge1$，取 $x_n\in B_{n,\overline K/K}$，使其在 $\mathcal O_{\overline K}/p^n$ 中的像等于 $\beta_n$ 的类。则 $y_n=(x_n)^{p^n}$ 只依赖于 $\beta_n$，与 $x_n$ 的选择无关，且
\[
 y:=(y_n)_n\in\varprojlim_n B_{n,\overline K/K}=B_{\infty,\overline K/K}.
\]
并且 $y\equiv1\pmod{J^{[1]}_{\infty,\overline K/K}}$。定义 $\varepsilon(\beta)$ 为 $y$ 在对数映射
\[
 1+J^{[1]}_{\infty,\overline K/K}\longrightarrow J^{[1]}_{\infty,\overline K/K}
\]
下的像。

\subsection*{§1.2. 对偶指数映射}
\textbf{1.2.1。} 对 $\Gal(\overline K/K)$ 的 de Rham 表示 $V$，典范映射
\[
 \exp:D_{\mathrm{dR}}(V)/D^0_{\mathrm{dR}}(V)\longrightarrow H^1(K,V),\qquad
 \exp^*:H^1(K,V)\longrightarrow D^0_{\mathrm{dR}}(V)
\]
分别称为\term{指数映射}{exponential map}{application exponentielle}与\term{对偶指数映射}{dual exponential map}{application exponentielle duale}。它们在 Galois 上同调与“de Rham 对象”之间起着重要的联系作用。§1.2 的目的是定义对偶指数映射。指数映射已在 [BK] 中定义，将在 §1.3 回顾；1.4.1(4) 将给出二者之间的对偶关系。

\textbf{1.2.2。} 令 $\chi_{\mathrm{cyclo}}:\Gal(\overline K/K)\longrightarrow\mathbb Z_p^\times$ 表示对所有 $n\ge1$ 的 $p^n$ 次单位根的作用，并令
\[
 \log(\chi_{\mathrm{cyclo}})\in H^1(K,\mathbb Z_p)=\Hom_{\mathrm{cont}}(\Gal(\overline K/K),\mathbb Z_p)
\]
为同态 $\sigma\mapsto\log(\chi_{\mathrm{cyclo}}(\sigma))$。
\begin{statement}{命题 1.2.3}
设 $V$ 是 $\Gal(\overline K/K)$ 的 de Rham 表示。则对所有 $i\in\mathbb Z$，与 $\log(\chi_{\mathrm{cyclo}})$ 的杯积给出同构
\[
 D^i_{\mathrm{dR}}(V)=H^0(K,V\otimes_{\mathbb Q_p}B^i_{\mathrm{dR}})\xrightarrow{\sim}H^1(K,V\otimes_{\mathbb Q_p}B^i_{\mathrm{dR}}).
\]
\end{statement}
\origpage{125}{120}
这里 $V\otimes B^i_{\mathrm{dR}}$ 赋予下文 1.2.5 中说明的拓扑。

\textbf{1.2.4。} 对 $\Gal(\overline K/K)$ 的 de Rham 表示 $V$，定义其对偶指数映射 $\exp^*$ 为复合
\[
 H^1(K,V)\longrightarrow H^1(K,V\otimes B^+_{\mathrm{dR}})\xrightarrow{\sim}D^0_{\mathrm{dR}}(V),
\]
第二个箭头是与 $\log(\chi_{\mathrm{cyclo}})$ 的杯积的逆（1.2.3）。

\textbf{1.2.5。} $V\otimes B^i_{\mathrm{dR}}$ 上的拓扑如下定义。先将 $B^+_{\mathrm{dR}}$ 看作拓扑环
$(B_{\infty,\overline K/K}/J^{[r]}_{\infty,\overline K/K})\otimes\mathbb Q$ 的逆极限来定义其拓扑。其中各商环张量 $\mathbb Q$ 后的拓扑，以
\[
 \{p^n(B_{\infty,\overline K/K}/J^{[r]}_{\infty,\overline K/K})\}_{n\ge0}
\]
为零点的基本邻域系。对 $K$ 的任意有限扩张 $L$，从 $K$ 的 $B^+_{\mathrm{dR}}$ 到 $L$ 的 $B^+_{\mathrm{dR}}$ 的双射是同胚。

对 $i\in\mathbb Z$，在 $B^i_{\mathrm{dR}}$ 上赋予如下拓扑：对 $B^+_{\mathrm{dR}}$-模 $B^i_{\mathrm{dR}}$ 的任一生成元 $a$，映射 $x\mapsto ax$ 从 $B^+_{\mathrm{dR}}$ 到 $B^i_{\mathrm{dR}}$ 都是同胚。（这一拓扑的存在来自以下事实：对任意 $r\ge0$，
\[
 J^{[r]}_{\infty,\overline K/K}/J^{[r+1]}_{\infty,\overline K/K}
\]
是 $\mathcal O_{C_p}=\varprojlim_n\mathcal O_{\overline K}/p^n\mathcal O_{\overline K}$ 上的可逆模（1.1.1）。因此对 $r,i\ge0$ 以及任一生成 $B^+_{\mathrm{dR}}$-模 $B^i_{\mathrm{dR}}$ 的 $a\in J^{[i]}_{\infty,\overline K/K}$，乘法映射
\[
 J^{[r]}_{\infty,\overline K/K}/J^{[r+1]}_{\infty,\overline K/K}
 \xrightarrow{x\mapsto ax}
 J^{[r+i]}_{\infty,\overline K/K}/J^{[r+i+1]}_{\infty,\overline K/K}
\]
是单射，其余核被 $p$ 的某个幂消去。当 $i\ge0$ 时，$B^i_{\mathrm{dR}}$ 的此拓扑等于它作为 $B^+_{\mathrm{dR}}$ 的子空间的拓扑。）

最后在 $V\otimes B^i_{\mathrm{dR}}$ 上赋予拓扑，使得对 $V$ 的任一 $\mathbb Q_p$-基 $(e_j)_j$，映射
\[
 \prod_j B^i_{\mathrm{dR}}\longrightarrow V\otimes B^i_{\mathrm{dR}},\qquad
 (x_j)_j\longmapsto\sum_j e_j\otimes x_j
\]
都是同胚。

\textbf{1.2.6。} 为证明命题 1.2.3，我们回顾 $B_{\mathrm{dR}}$ 和 de Rham 表示的一些事实。
\origpage{126}{121}
令 $C_p$ 为 $\overline K$ 的完备化。存在 $\Gal(\overline K/K)$-模的典范 $C_p$-线性同构
\[
 C_p(i)\simeq B^i_{\mathrm{dR}}/B^{i+1}_{\mathrm{dR}}\qquad(i\in\mathbb Z),
\]
其定义如下。$i=0$ 时，它由同构 $B_{n,\overline K/K}/J_{n,\overline K/K}\xrightarrow{\sim}\mathcal O_{\overline K}/p^n$ 诱导；$i=1$ 时，由映射 $\varepsilon:\mathbb Q_p(1)\longrightarrow B^1_{\mathrm{dR}}/B^2_{\mathrm{dR}}$ 诱导；一般的 $i$ 则利用 $B_{\mathrm{dR}}$ 中的乘法，由 $i=0,1$ 的情形诱导。对 de Rham 表示 $V$，取保持滤过的同构
$V\otimes_{\mathbb Q_p}B_{\mathrm{dR}}\simeq D_{\mathrm{dR}}(V)\otimes_K B_{\mathrm{dR}}$（1.1.3）的分次商，得到\term{Hodge--Tate 分解}{Hodge--Tate decomposition}{décomposition de Hodge--Tate}
\begin{equation*}
 V\otimes_{\mathbb Q_p}(B^i_{\mathrm{dR}}/B^{i+1}_{\mathrm{dR}})
 \simeq\bigoplus_{r+s=i}\operatorname{gr}^r(D_{\mathrm{dR}}(V))\otimes_K C_p(s).\tag{1}
\end{equation*}
\textbf{1.2.7。} 下面证明命题 1.2.3。关键是 Tate [Ta$_1$] 的如下结果：当 $(m,r)\ne(0,0),(1,0)$ 时，连续上同调 $H^m(K,C_p(r))$ 为零；而 $H^0(K,C_p)$ 与 $H^1(K,C_p)$ 都是一维 $K$-向量空间，分别由 $1$ 和 $\log(\chi_{\mathrm{cyclo}})$ 生成。

由此及 1.2.6(1)，得到正合列
\[
 0\longrightarrow H^1(K,V\otimes B^{i+1}_{\mathrm{dR}})
 \longrightarrow H^1(K,V\otimes B^i_{\mathrm{dR}})
 \longrightarrow D^i_{\mathrm{dR}}(V)/D^{i+1}_{\mathrm{dR}}(V)\longrightarrow0.
\]
因此当 $i$ 充分小时，\footnote{译注：原书此处写“$i$ 充分小”，且随后的上同调次数与极限均如下式；原书如此，照录。}
\[
 H^1(K,V\otimes B^i_{\mathrm{dR}})\simeq\varprojlim_n H^1(K,V\otimes B^n_{\mathrm{dR}})=0,
\]
再由这一正合列即得 1.2.3。

\subsection*{§1.3. 指数映射}
我们回顾 Fontaine 的环 $B_{\mathrm{crys}}$ 与函子 $D_{\mathrm{crys}}$（[Fo$_1$]、[FM]），以及 de Rham 表示的指数映射（[BK]）。

\origpage{127}{122}
\textbf{1.3.1。} $B_{\mathrm{dR}}$ 的子环 $B_{\mathrm{crys}}$ 如下定义。考虑环 $B_{\infty,\overline K/K_0}$（注意 $\overline K$ 也是 $K_0$ 的代数闭包）。映射
$B_{\infty,\overline K/K_0}\longrightarrow B_{\infty,\overline K/K}$ 为单射，且 $\varepsilon:\mathbb Z_p(1)\longrightarrow B_{\infty,\overline K/K}$ 的像包含于 $B_{\infty,\overline K/K_0}$。定义
\[
 B_{\mathrm{crys}}=B_{\infty,\overline K/K_0}[p^{-1},t^{-1}]\subset B_{\mathrm{dR}},
\]
其中 $t$ 为 $\mathbb Z_p(1)$ 的任一基在 $\varepsilon$ 下的像；这样 $B_{\mathrm{crys}}$ 与 $t$ 的选择无关。

Frobenius 算子 $\varphi:B_{\mathrm{crys}}\longrightarrow B_{\mathrm{crys}}$ 由
$\mathcal O_{\overline K}/p\longrightarrow\mathcal O_{\overline K}/p,\ x\mapsto x^p$、$W_n(k)$ 的 Frobenius，以及同构
\begin{align*}
 &H^0(\Spec(\mathcal O_{\overline K}/p^n)/\Spec(W_n(k))_{\mathrm{crys}},\mathcal O_{\mathrm{crys}})\\
 &\qquad\xrightarrow{\sim}H^0(\Spec(\mathcal O_{\overline K}/p)/\Spec(W_n(k))_{\mathrm{crys}},\mathcal O_{\mathrm{crys}})
\end{align*}
诱导。

\textbf{1.3.2。} 对带有 $\Gal(\overline K/K)$ 连续作用的有限维 $\mathbb Q_p$-向量空间 $V$，令
\[
 D_{\mathrm{crys}}(V)=H^0(K,V\otimes_{\mathbb Q_p}B_{\mathrm{crys}}).
\]
这是一个有限维 $K_0$-向量空间，带有 Frobenius 线性算子 $1\otimes\varphi$，仍记为 $\varphi$。“Frobenius 线性”指对所有 $a\in K_0$ 和 $x\in D_{\mathrm{crys}}(V)$，都有 $\varphi(ax)=\varphi(a)\varphi(x)$。

\textbf{1.3.3。} 序列
\begin{equation*}
 0\longrightarrow\mathbb Q_p\xrightarrow{\alpha}
 B_{\mathrm{crys}}\oplus B^+_{\mathrm{dR}}\xrightarrow{\beta}
 B_{\mathrm{crys}}\oplus B_{\mathrm{dR}}\longrightarrow0\tag{1}
\end{equation*}
正合，其中 $\alpha(x)=(x,x)$，$\beta(x,y)=((1-\varphi)(x),x-y)$。在 [BK] 中，(1) 的正合性利用 Fontaine 和 Messing [FM] 的关键正合列
\begin{equation*}
 0\longrightarrow\mathbb Q_p(r)\longrightarrow J^{[r]}_{\infty,\overline K/K_0}\otimes\mathbb Q
 \xrightarrow{1-p^{-r}\varphi}B_{\infty,\overline K/K_0}\otimes\mathbb Q\longrightarrow0
 \quad(r\ge0)\tag{2}
\end{equation*}
得到证明。[FM] 未写出 (2) 正合性的证明，但 de Shalit 最近在 [dS$_3$] 中写出了证明。
\origpage{128}{123}
对带有 $\Gal(\overline K/K)$ 连续作用的有限维 $\mathbb Q_p$-向量空间 $V$，将正合列 (1) 与 $V$ 在 $\mathbb Q_p$ 上张量，得到正合列
\begin{equation*}
\begin{aligned}
 0\longrightarrow H^0(K,V)&\xrightarrow{\alpha}D_{\mathrm{crys}}(V)\oplus D^0_{\mathrm{dR}}(V)\\
 &\xrightarrow{\beta}D_{\mathrm{crys}}(V)\oplus D_{\mathrm{dR}}(V)\xrightarrow{\delta}H^1(K,V).
\end{aligned}\tag{3}
\end{equation*}
\textbf{1.3.4。} 对 1.3.3 中的 $V$，定义指数映射
\[
 \exp:D_{\mathrm{dR}}(V)/D^0_{\mathrm{dR}}(V)\longrightarrow H^1(K,V)
\]
为 $\exp(x\bmod D^0_{\mathrm{dR}}(V))=\delta(0,x)$，其中 $x\in D_{\mathrm{dR}}(V)$。

\textbf{例 1.3.5。} 设 $A$ 是 $K$ 上的阿贝尔簇，令 $V_pA=T_pA\otimes\mathbb Q$，其中 $T_pA$ 是 $A$ 的 $p$ 进\term{Tate 模}{Tate module}{module de Tate}。则 $V_pA$ 是 $\Gal(\overline K/K)$ 的 de Rham 表示，$\operatorname{Lie}(A)$ 与 $D_{\mathrm{dR}}(V_pA)/D^0_{\mathrm{dR}}(V_pA)$ 等同（[Fo$_1$]），而
\[
 \exp:D_{\mathrm{dR}}(V_pA)/D^0_{\mathrm{dR}}(V_pA)\longrightarrow H^1(K,V_pA)
\]
与复合映射
\[
 \operatorname{Lie}(A)\longrightarrow A(K)\otimes\mathbb Q\longrightarrow H^1(K,V_pA)
\]
等同。其中第一个箭头是 $p$ 进 Lie 群 $A(K)$ 的指数映射，第二个来自 $A$ 的 Kummer 序列（参见第一章 §2.3）。

\textbf{例 1.3.6。} 设 $G$ 是 $\mathcal O_K$ 上的一个\term{$p$ 可除群}{$p$-divisible group}{groupe $p$-divisible}，$T_pG$ 是其 Tate 模，$V_pG=T_pG\otimes\mathbb Q$。则 $V_pG$ 是 $\Gal(\overline K/K)$ 的 de Rham 表示，$\operatorname{Lie}(G)$ 与 $D_{\mathrm{dR}}(V_pG)/D^0_{\mathrm{dR}}(V_pG)$ 等同（[Fo$_1$]），指数映射与
\[
 \operatorname{Lie}(G)\longrightarrow\bigl(\varprojlim_n G(\mathcal O_K/p^n)\bigr)\otimes\mathbb Q
 \longrightarrow H^1(K,V_pG)
\]
等同。第一个箭头是 $\mathcal O_K$ 上与 $G$ 的连通部分相伴的形式群的指数映射（[Ta$_2$] §2.2），第二个来自 $G$ 的 Kummer 序列
\[
 0\longrightarrow G_n\longrightarrow G\xrightarrow{p^n}G\longrightarrow0.
\]
关于这两个例子中的指数映射，参见 [BK] §3。

\origpage{129}{124}
\subsection*{§1.4. Tate 对偶}
本节假设 $k$ 是有限域。以下 1.4.1 的 (1)--(3) 是 Tate 的对偶定理（参见 [Se] 第二章），(4) 则将在本节证明。
\begin{statement}{定理 1.4.1}
设 $\ell$ 为素数。
\begin{enumerate}[label=(\arabic*)]
\item 存在典范同构 $H^2(K,\mathbb Q_\ell(1))\simeq\mathbb Q_\ell$，其刻画见下文 1.4.2。
\item 设 $V$ 是带有 $\Gal(\overline K/K)$ 连续作用的有限维 $\mathbb Q_\ell$-向量空间。则对所有 $q$，$H^q(K,V)$ 都是有限维 $\mathbb Q_\ell$-向量空间，且当 $q\ne0,1,2$ 时为零。
\item 令 $V$ 如 (2)，$V^*=\Hom_{\mathbb Q_\ell}(V,\mathbb Q_\ell)$ 为其对偶表示，$\sigma\in\Gal(\overline K/K)$ 的作用为 $h\mapsto h\circ\sigma^{-1}$。则杯积
\[
 H^q(K,V)\times H^{2-q}(K,V^*(1))\longrightarrow H^2(K,\mathbb Q_\ell(1))\simeq\mathbb Q_\ell
\]
为\term{完美配对}{perfect pairing}{accouplement parfait}。
\item 假设 $\ell=p$，且 (2) 中的 $V$ 是 de Rham 表示。则 $\exp^*:H^1(K,V)\longrightarrow D_{\mathrm{dR}}(V)$ 等于复合映射
\begin{align*}
 H^1(K,V)&\xrightarrow{\sim}\Hom_{\mathbb Q_p}(H^1(K,V^*(1)),\mathbb Q_p)\\
 &\longrightarrow\Hom_{\mathbb Q_p}(D_{\mathrm{dR}}(V^*(1)),\mathbb Q_p)\xrightarrow{\sim}D_{\mathrm{dR}}(V).
\end{align*}
第一个箭头来自 (3) 的对偶，第二个是 $V^*(1)$ 的指数映射的 $\mathbb Q_p$-对偶，第三个由典范配对
\[
 D_{\mathrm{dR}}(V)\times D_{\mathrm{dR}}(V^*(1))\longrightarrow D_{\mathrm{dR}}(\mathbb Q_p(1))\simeq K\xrightarrow{\mathrm{trace}}\mathbb Q_p
\]
诱导。
\end{enumerate}
\end{statement}

\origpage{130}{125}
\textbf{1.4.2。} 1.4.1(1) 的同构由以下性质刻画。设 $\pi$ 为 $K$ 的素元，$\{\pi\}\in H^1(K,\mathbb Q_\ell(1))$ 为 $\pi$ 在 Kummer 序列
\[
 0\longrightarrow\mathbb Z/\ell^n(1)\longrightarrow\mathbb G_m\xrightarrow{\ell^n}\mathbb G_m\longrightarrow0
\]
的连接映射诱导的 $K^\times\longrightarrow H^1(K,\mathbb Q_\ell(1))$ 下的像。又设
\[
 \chi\in H^1(K,\mathbb Q_\ell)=\Hom_{\mathrm{cont}}(\Gal(\overline K/K),\mathbb Q_\ell)
\]
是经由典范满射 $\Gal(\overline K/K)\longrightarrow\Gal(K_{\mathrm{ur}}/K)$ 分解的同态。那么 1.4.1(1) 的同构将杯积 $\{\pi\}\cup\chi$ 送到 $\chi(\operatorname{Frob}_K)\in\mathbb Q_\ell$；这里 $\operatorname{Frob}_K\in\Gal(K_{\mathrm{ur}}/K)$ 是算术 Frobenius，即在 $K_{\mathrm{ur}}$ 的剩余域 $\overline k$ 上诱导 $x\mapsto x^{\#k}$ 的唯一元素。

现在证明 1.4.1(4)。
\begin{statement}{引理 1.4.3}
设 $\delta:H^1(K,B^+_{\mathrm{dR}}/B^1_{\mathrm{dR}})\longrightarrow H^2(K,\mathbb Q_p(1))$ 为正合列
\[
 0\longrightarrow\mathbb Q_p(1)\longrightarrow B_{\mathrm{crys}}^{\varphi=p}\cap B^+_{\mathrm{dR}}
 \longrightarrow B^+_{\mathrm{dR}}/B^1_{\mathrm{dR}}\longrightarrow0
\]
的连接同态。则对 $\Gal(\overline K/K)$ 的任意 de Rham 表示 $V$，任取 $a\in D_{\mathrm{dR}}(V)$ 和 $b\in H^1(K,V^*(1))$，在 $H^2(K,\mathbb Q_p(1))$ 中都有
\[
 \exp(a)\cup b=\delta(\langle a,\exp^*(b)\rangle\log(\chi_{\mathrm{cyclo}})).
\]
这里 $\langle\ ,\ \rangle$ 是配对
$D_{\mathrm{dR}}(V)\times D_{\mathrm{dR}}(V^*(1))\longrightarrow D_{\mathrm{dR}}(\mathbb Q_p(1))=K$。
\end{statement}
\begin{proof}
取 $i<0$ 使 $a\in D^i_{\mathrm{dR}}(V)$。由上同调的形式运算，$\exp(a)\cup b$ 等于 $b$ 在下列复合映射中的像：
\begin{align*}
 H^1(K,V^*(1))&\longrightarrow H^1(K,V^*(1)\otimes B^+_{\mathrm{dR}})\\
 &\xrightarrow{a\cup}H^1(K,\mathbb Q_p(1)\otimes B^i_{\mathrm{dR}})
 =H^1(K,B^{i+1}_{\mathrm{dR}})\\
 &\longrightarrow H^1(K,B^{i+1}_{\mathrm{dR}}/B^1_{\mathrm{dR}})
 \xleftarrow{\sim}H^1(K,B^+_{\mathrm{dR}}/B^1_{\mathrm{dR}})
 \xrightarrow{\delta}H^2(K,\mathbb Q_p(1)).
\end{align*}
其中 $a\cup$ 是与 $a$ 的杯积。因为 $b$ 在 $H^1(K,V^*(1)\otimes B^+_{\mathrm{dR}})$ 中的像为 $\exp^*(b)\log(\chi_{\mathrm{cyclo}})$，故得引理。
\end{proof}

\origpage{131}{126}
\textbf{1.4.4。} 由 1.4.3，为证明 1.4.1(4)，只须证明复合
\[
 H^1(K,B^+_{\mathrm{dR}}/B^1_{\mathrm{dR}})\xrightarrow{\delta}H^2(K,\mathbb Q_p(1))\xrightarrow{\sim}\mathbb Q_p
\]
将 $a\log(\chi_{\mathrm{cyclo}})$ 送到 $-\Tr_{K/\mathbb Q_p}(a)$，其中 $a\in K$。负号来自
\[
 x\cup y=-y\cup x\quad\text{于 }H^2(K,\mathbb Q_p(1)),\qquad
 x\in H^1(K,V),\quad y\in H^1(K,V^*(1)).
\]
在 1.4.3 中取 $V=\mathbb Q_p(1)$、$b=\log(\chi_{\mathrm{cyclo}})$，可得
\[
 \delta(a\log(\chi_{\mathrm{cyclo}}))=\exp(a)\cup\log(\chi_{\mathrm{cyclo}})
 \quad\text{于 }H^2(K,\mathbb Q_p(1)),
\]
其中 $a\in K$，$\exp$ 是指数映射 $K\longrightarrow H^1(K,\mathbb Q_p(1))$。因此归结为：
\begin{statement}{引理 1.4.5}
对于 $V=\mathbb Q_p(1)$，1.4.1(3) 中的配对
\[
 H^1(K,\mathbb Q_p(1))\times H^1(K,\mathbb Q_p)\longrightarrow H^2(K,\mathbb Q_p(1))\simeq\mathbb Q_p
\]
将 $(\exp(a),\log(\chi_{\mathrm{cyclo}}))$ 送到 $-\Tr_{K/\mathbb Q_p}(a)$，对所有 $a\in K$ 成立。
\end{statement}
\begin{proof}
利用 $K/\mathbb Q_p$ 的迹映射，可归结到 $K=\mathbb Q_p$。在此情形下，结论是 $\mathbb Q_p$ 的\term{局部类域论}{local class field theory}{théorie locale du corps de classes}中的熟知结果。
\end{proof}

\subsection*{§2. Lubin--Tate 群的广义显式互反律}
\addcontentsline{toc}{subsection}{§2. Lubin--Tate 群的广义显式互反律}
\term{显式互反律}{explicit reciprocity law}{loi de réciprocité explicite}是局部域算术中的一个重要课题，已有许多人的研究。经典形式是一条利用微分形式显式描述 Hilbert 符号的公式。我们的观点是：Hilbert 符号与微分形式之间这种神秘的关系，与 $p$ 进周期理论中 $p$ 进平展上同调和微分形式之间的神秘关系密切相关。
\origpage{132}{127}
这一观点在 [Ka$_1$] 中得到强调。注意，Hilbert 符号与 Galois 上同调理论密切相关，而后者是平展上同调理论的一个分支。

我相信 $\Gal(\overline K/K)$ 的所有 $p$ 进表示都存在显式互反律，尽管我还不能将它们表述出来。对于 de Rham 表示 $V$，这一互反律应当是对 $D_{\mathrm{dR}}(V)$ 与 $V$ 的 Galois 上同调之间关系的某种显式描述；更确切地说，应是 $V$ 的映射 $\exp$ 和 $\exp^*$ 的某种显式描述。本节将对
$T^{\otimes(-r)}(1)\otimes\mathbb Q$（$r\ge1$）的 $\exp^*$ 给出显式描述（定理 2.1.7），其中 $T$ 是一个\term{Lubin--Tate 群}{Lubin--Tate group}{groupe de Lubin--Tate}的 Tate 模。定理 2.1.7 在 $r=1$ 时包含 Wiles [Wi$_1$] 的显式互反律，也包含 [BK] §2 的显式互反律；后者研究的 Lubin--Tate 群是形式乘法群，并且考虑的是 $\mathbb Q_p$ 的非分歧扩张的 Galois 上同调。

de Shalit [dS$_2$] 第一章，是关于 Lubin--Tate 群及 Wiles 显式互反律的极好入门读物。关于显式互反律，已有非常多的工作。遗憾的是，这里未能将它们一一列出；其中一些可见 [dS$_1$]、[dS$_2$]、[Ka$_1$]。

\subsection*{§2.1. 结果的陈述}
这里不再回顾 Lubin--Tate 群的定义，其基本事实参见 [dS$_2$]。

\textbf{2.1.1。} 设 $F$ 为 $K$ 的子域，满足 $[F:\mathbb Q_p]<\infty$，并假设 $F$ 的素元也是 $K$ 的素元。固定 $F$ 的一个素元 $\pi$。令 $\Gamma$ 为对应于 $(F,\pi)$ 的 $\mathcal O_F$ 上的 Lubin--Tate 群。

\origpage{133}{128}
令 $G$ 为经标量扩张 $\mathcal O_F\longrightarrow\mathcal O_K$ 从 $\Gamma$ 得到的 $\mathcal O_K$ 上的连通 $p$ 可除群。将 $a\in\mathcal O_F$ 对 $G$ 的作用记为 $[a]:G\longrightarrow G$。

\textbf{2.1.2。} 令 $T$ 为 $G$ 的 Tate 模，$V=T\otimes\mathbb Q$。对 $m\in\mathbb Z$，以 $T^{\otimes m}$、$V^{\otimes m}$、$\operatorname{coLie}(G)^{\otimes m}$ 分别表示 $T$、$V$、$\operatorname{coLie}(G)$ 作为可逆的 $\mathcal O_F$-模、$F$-模、$\mathcal O_K$-模的第 $m$ 个张量幂。对 $n\ge0$，令 $K_n$ 为对应于
\[
 \Ker\bigl(\Gal(\overline K/K)\longrightarrow\operatorname{Aut}(T/\pi^nT)\bigr)
\]
的 $K$ 的扩张。下文用 $\xi$ 表示 $T$ 的一个 $\mathcal O_F$-基。对 $n\ge0$，令
\[
 \xi_n=\pi^{-n}\xi\bmod T\in V/T=G(\overline K).
\]
\textbf{2.1.3。} 对 $r\ge1$，如后文 2.2.3 将说明的那样，有典范同构
\begin{equation*}
 D^0_{\mathrm{dR}}(V^{\otimes(-r)}(1))\simeq\operatorname{coLie}(G)^{\otimes r}\otimes_{\mathcal O_K}K.\tag{1}
\end{equation*}
它对每个 $n\ge0$ 都诱导同构
\begin{equation*}
 D^0_{\mathrm{dR}}(K_n,V^{\otimes(-r)}(1))\simeq\operatorname{coLie}(G)^{\otimes r}\otimes_{\mathcal O_K}K_n.\tag{2}
\end{equation*}
这里 $D_{\mathrm{dR}}(K_n,\ )$ 表示用于 $\Gal(\overline K/K_n)$ 的表示的 $D_{\mathrm{dR}}$，参见 1.1.3。

\textbf{2.1.4。} 用 $\mathcal O(G)$ 表示 $G$ 上的函数环，即有限群概形 $\Ker(p^n:G\longrightarrow G)$ 的仿射坐标环的逆极限，因此 $\mathcal O(G)\simeq\mathcal O_K[[t]]$。用 $\phi:\mathcal O(G)\longrightarrow\mathcal O(G)$ 表示唯一的环同态，它在 $\mathcal O(\Gamma)$ 上为恒等，且对任意 $f\in\mathcal O(G)$，满足
\[
 \phi(f)\equiv f^q\pmod{\pi\mathcal O(G)},
\]
其中 $q$ 为 $F$ 的剩余域的阶。\translatornote{11}

\textbf{2.1.5。} 对关于范数相容的系统 $u=(u_n)_n\in\varprojlim_n K_n^\times$，定义
\[
 s_{r,n}(u,\xi)\in H^1(K_n,T^{\otimes(-r)}(1))\qquad(n\ge0,\ r\in\mathbb Z)
\]
为 $u$ 在以下映射中的像：
\origpage{134}{129}
\begin{align*}
 \varprojlim_m K_m^\times
 &\longrightarrow\varprojlim_m H^1(K_m,\mathbb Z_p(1))\\
 &\xrightarrow{(*)}\varprojlim_m H^1(K_m,T^{\otimes(-r)}(1)/\pi^m)\\
 &\xrightarrow{\mathrm{trace}}\varprojlim_m H^1(K_n,T^{\otimes(-r)}(1)/\pi^m)
 \xleftarrow{\sim}H^1(K_n,T^{\otimes(-r)}(1)).
\end{align*}
前三个逆极限分别关于范数或迹取得，而 $(*)$ 由与 $(\xi_m)^{\otimes(-r)}$ 的杯积定义。

\textbf{2.1.6。} 对上述 $u$，用 $g_{u,\xi}$ 表示 $u$ 关于 $\xi$ 的\term{Coleman 幂级数}{Coleman power series}{série de Coleman}。它是 $\mathcal O(G)[1/h]$ 的一个可逆元，其中 $h$ 为\term{增广理想}{augmentation ideal}{idéal d'augmentation} $\Ker(\mathcal O(G)\longrightarrow\mathcal O_K)$ 的生成元，并由以下性质刻画：
\[
 (\phi^{-n}(g_{u,\xi}))(\xi_n)=u_n\qquad(n\ge1).
\]
参见 [Co]。

现在给出本节的主要定理。
\begin{statement}{定理 2.1.7（显式互反律）}
对 $n,r\ge1$，
\begin{align*}
 \exp^*:H^1(K_n,T^{\otimes(-r)}(1))
 &\longrightarrow D^0_{\mathrm{dR}}(K_n,V^{\otimes(-r)}(1))\\
 &\simeq\operatorname{coLie}(G)^{\otimes r}\otimes_{\mathcal O_K}K_n
\end{align*}
将 $s_{r,n}(u,\xi)$ 送到
\[
 \frac{1}{(r-1)!}\pi^{-nr}\omega^{\otimes r}\otimes
 \left(\left(\frac d\omega\right)^r\log(\phi^{-n}(g_{u,\xi}))\right)(\xi_n).
\]
这里 $\omega$ 是 $\operatorname{coLie}(G)=\Hom_{\mathcal O_K}(\operatorname{Lie}(G),\mathcal O_K)$ 的任一 $\mathcal O_K$-基，且
\[
 \left(\frac d\omega\right)^r\log(\ )=
 \left(\frac d\omega\right)^{r-1}\left(\frac{d\log(\ )}{\omega}\right),
\]
其中 $(d/\omega)^{r-1}$ 是算子 $d/\omega:f\mapsto df/\omega$ 迭代 $r-1$ 次。此比值通过将 $\omega$ 看作 $G$ 上的不变微分形式取得。
\end{statement}
定理 2.1.7 的 $r=1$ 情形就是 Wiles 的显式互反律（[Wi$_1$]）。

\subsection*{§2.2. Lubin--Tate 群与晶体上同调}
我们利用 [Ta$_2$]、[Fo$_1$]、[Fo$_2$] 中的理论定义同构 2.1.3(1)，然后陈述定理 2.2.7，并说明如何由它推出定理 2.1.7。

\origpage{135}{130}
\textbf{2.2.1。} 令 $\Omega^1(G)$ 为 $G$ 上微分形式的空间。注意
\[
 \Omega^1(G)\simeq\mathcal O(G)\otimes_{\mathcal O_K}\operatorname{coLie}(G).
\]
对 $n\ge0$，令 $S_n$ 为一个 $p$ 进完备且分离的拓扑环，由以下性质刻画：对于每个 $i\ge1$，$\Spec(S_n/p^iS_n)$ 是
$\Spec(\mathcal O_{K_n}/p^i\mathcal O_{K_n})$ 在 $\Spec(\mathcal O(G)/p^i\mathcal O(G))$ 中，关于由 $\xi_n:\Spec(\mathcal O_{K_n})\longrightarrow G$ 诱导的嵌入的\term{分幂包络}{PD envelope}{enveloppe à puissances divisées}。

对 $r\ge1$，令 $J^{[r]}_{S_n}\subset S_n$ 为 $\Ker(S_n/p^iS_n\longrightarrow\mathcal O_{K_n}/p^i\mathcal O_{K_n})$ 的第 $r$ 个分幂关于 $i$ 的逆极限。由于 $G$ 在 $\mathcal O_K$ 上形式光滑，晶体上同调的基本定理给出
\begin{align*}
 &R\Gamma(\Spec(\mathcal O_{K_n}/\pi^i)/\Spec(\mathcal O_K/\pi^i)_{\mathrm{crys}},J^{[r]})\\
 &\qquad\simeq\left[J^{[r]}_{S_n}/\pi^i\xrightarrow d
 J^{[r-1]}_{S_n}/\pi^i\otimes_{\mathcal O(G)}\Omega^1(G)\right]\qquad(i\ge1),
\end{align*}
其中 $J=\Ker(\mathcal O_{\mathrm{crys}}\longrightarrow\mathcal O_{\mathrm{zar}})$。

\textbf{2.2.2。} 对 $\omega\in\operatorname{coLie}(G)$，令 $\ell_\omega\in J_{S_0}$ 为 $G$ 关于 $\omega$ 的对数。于是 $d\ell_\omega=\omega$。对 $n\ge0$，令
\[
 \ell_{\omega,\xi,n}=[\pi^n]^*\ell_\omega\in J_{S_n}.
\]
则 $d\ell_{\omega,\xi,n}=\pi^n\omega$；特别地，$d\ell_{\omega,\xi,n}\equiv0\pmod{\pi^n}$。因此由 2.2.1，$\ell_{\omega,\xi,n}\bmod\pi^n$ 可看作
\[
 H^0(\Spec(\mathcal O_{K_n}/\pi^n)/\Spec(\mathcal O_K/\pi^n)_{\mathrm{crys}},J)
\]
中的元素。它在
\begin{align*}
 &H^0(\Spec(\mathcal O_{\overline K}/\pi^n)/\Spec(\mathcal O_K/\pi^n)_{\mathrm{crys}},J)\\
 &\qquad=J_{n,\overline K/K}=J_{\infty,\overline K/K}/\pi^n
\end{align*}
中的像，仍记为 $\ell_{\omega,\xi,n}$。

\origpage{136}{131}
\textbf{2.2.3。} 有配对
\begin{equation*}
 \operatorname{coLie}(G)\otimes_{\mathcal O_F}T\longrightarrow J_{\infty,\overline K/K},\qquad
 \omega\otimes\xi\longmapsto(\ell_{\omega,\xi,n})_n,\tag{1}
\end{equation*}
诱导映射
\begin{equation*}
 \operatorname{coLie}(G)\longrightarrow T^{\otimes(-1)}\otimes_{\mathcal O_F}J_{\infty,\overline K/K}.\tag{2}
\end{equation*}
由 [Ta$_2$]、[Fo$_1$]、[Fo$_2$]，有以下 (3) 与 (4)：
\begin{enumerate}[label=(\arabic*),start=3]
\item 映射 (2) 诱导可逆 $B^+_{\mathrm{dR}}$-模的同构
\[
 \operatorname{coLie}(G)\otimes_{\mathcal O_K}B^+_{\mathrm{dR}}\xrightarrow{\sim}
 V^{\otimes(-1)}\otimes_F B^1_{\mathrm{dR}}.
\]
\item $V^{\otimes(-1)}\otimes_{\mathbb Q_p}B^+_{\mathrm{dR}}\longrightarrow V^{\otimes(-1)}\otimes_F B^+_{\mathrm{dR}}$ 的核，作为带有 $\Gal(\overline K/K)$ 作用的 $B^+_{\mathrm{dR}}$-模，同构于有限多个 $B^+_{\mathrm{dR}}$ 的直和。
\end{enumerate}
分别在 $B^+_{\mathrm{dR}}$ 和 $F\otimes_{\mathbb Q_p}B^+_{\mathrm{dR}}$ 上取第 $r$ 个张量幂，由 (3)、(4) 分别得到以下 (5)、(6)：
\begin{enumerate}[label=(\arabic*),start=5]
\item $\operatorname{coLie}(G)^{\otimes r}\otimes_{\mathcal O_K}B^+_{\mathrm{dR}}\xrightarrow{\sim}V^{\otimes(-r)}\otimes_F B^r_{\mathrm{dR}}$；
\item $V^{\otimes(-r)}\otimes_{\mathbb Q_p}B^+_{\mathrm{dR}}\longrightarrow V^{\otimes(-r)}\otimes_F B^+_{\mathrm{dR}}$ 的核，作为带有 $\Gal(\overline K/K)$ 作用的 $B^+_{\mathrm{dR}}$-模，同构于有限多个 $B^+_{\mathrm{dR}}$ 的直和。
\end{enumerate}
由 (5)、(6)，取 $H^0(K,\ )$，得到
\begin{align*}
 \operatorname{coLie}(G)^{\otimes r}\otimes_{\mathcal O_K}K
 &\xrightarrow{\sim}H^0(K,V^{\otimes(-r)}\otimes_F B^r_{\mathrm{dR}})\\
 &=H^0(K,V^{\otimes(-r)}\otimes_F B^1_{\mathrm{dR}})\\
 &=H^0(K,V^{\otimes(-r)}(1)\otimes_F B^+_{\mathrm{dR}})\\
 &\xleftarrow{\sim}H^0(K,V^{\otimes(-r)}(1)\otimes_{\mathbb Q_p}B^+_{\mathrm{dR}})
 =D^0_{\mathrm{dR}}(V^{\otimes(-r)}(1)),
\end{align*}
这就给出了 2.1.3(1) 中的同构。

\textbf{2.2.4。} 我们将需要 2.2.3 中某些向量空间的整结构。事实上，[Ta$_2$] 中的论证还证明了 1.2.7 所引 Tate 结果的以下整系数版本。
\origpage{137}{132}
设 $L$ 是特征零的完备离散赋值域，其剩余域是特征 $p>0$ 的完美域。令 $q\ge0$、$r\in\mathbb Z$。则：
\begin{enumerate}[label=(\arabic*)]
\item 若 $q\ge2$ 或 $r\ne0$，则 $H^q(L,\widehat{\mathcal O}_{\overline L}(r))$ 被 $p$ 的某个幂消去。
\item 若 $q\ge2$ 或 $r\ne0$，则存在整数 $m\ge0$，使得 $p^m$ 对所有 $n\ge0$ 都消去 $H^q(L,(\mathcal O_{\overline L}/p^n)(r))$。
\item $\mathcal O_L\xrightarrow{\sim}H^0(L,\widehat{\mathcal O}_{\overline L})$。
\item 映射 $\mathcal O_L\longrightarrow H^1(L,\widehat{\mathcal O}_{\overline L})$，$a\mapsto a\log(\chi_{\mathrm{cyclo}})$，为单射，其余核被 $p$ 的某个幂消去。
\end{enumerate}
由这些事实，对于 $q=0,1$、$r\ge1$ 及 $K$ 的任意有限扩张 $L$，以下映射的核与余核均被 $p$ 的某个幂消去：
\begin{align*}
 \operatorname{coLie}(G)^{\otimes r}\otimes_{\mathcal O_K}\mathcal O_L
 &\longrightarrow H^q\bigl(L,T^{\otimes(-r)}\otimes_{\mathcal O_F}
 J^{[r]}_{\infty,\overline K/K}/J^{[r+1]}_{\infty,\overline K/K}\bigr)\\
 &\longrightarrow H^q\bigl(L,T^{\otimes(-r)}\otimes_{\mathcal O_F}
 J_{\infty,\overline K/K}/J^{[r+1]}_{\infty,\overline K/K}\bigr).
\end{align*}
这里 $q=0$ 时，第一个箭头是由 2.2.3(2) 诱导的单射；$q=1$ 时，它是 $q=0$ 的映射与“同 $\log(\chi_{\mathrm{cyclo}})$ 作杯积”的复合。

\textbf{2.2.5。} 令
\[
 \Theta=\varprojlim(\cdots\longrightarrow\Omega^1(G)\longrightarrow\Omega^1(G)\longrightarrow\Omega^1(G)),
\]
其中箭头为乘法 $[\pi]:G\longrightarrow G$ 所伴随的迹映射 $\Tr_\pi:\Omega^1(G)\longrightarrow\Omega^1(G)$。因此在 $\Omega^1(G)$ 上有 $\Tr_\pi\circ[\pi]^*=\pi$。\translatornote{12}

\textbf{2.2.6。} 我们如下定义典范映射（$n\ge0$、$r\in\mathbb Z$）：
\begin{align*}
 \mu_{\xi,n}:S_n\otimes_{\mathcal O(G)}\Omega^1(G)
 &\longrightarrow H^1(K_n,J_{\infty,\overline K/K}),\\
 \mu_{r,\xi,n}:S_n\otimes_{\mathcal O(G)}\Omega^1(G)
 &\longrightarrow H^1\bigl(K_n,(T^{\otimes(-r)}\otimes_{\mathcal O_F}
 J_{\infty,\overline K/K}/J^{[r+1]}_{\infty,\overline K/K})/\pi^n\bigr),\\
 \theta_{r,\xi,n}:\Theta
 &\longrightarrow H^1\bigl(K_n,T^{\otimes(-r)}\otimes_{\mathcal O_F}
 J_{\infty,\overline K/K}/J^{[r+1]}_{\infty,\overline K/K}\bigr).
\end{align*}
\origpage{138}{133}
令
\begin{align*}
 S_n\otimes_{\mathcal O(G)}\Omega^1(G)&\longrightarrow
 H^1(\Spec(\mathcal O_{K_n}/\pi^i)/\Spec(\mathcal O_K/\pi^i)_{\mathrm{crys}},J)
 \quad(i\ge1)
\end{align*}
为由 2.2.1 的同构定义的映射。还有映射
\begin{align*}
 &R\Gamma(\Spec(\mathcal O_{K_n}/\pi^i)/\Spec(\mathcal O_K/\pi^i)_{\mathrm{crys}},J)\\
 &\quad\longrightarrow R\Gamma\bigl(\Gal(\overline K/K_n),
 R\Gamma(\Spec(\mathcal O_{\overline K}/\pi^i)/\Spec(\mathcal O_K/\pi^i)_{\mathrm{crys}},J)\bigr)\\
 &\quad=R\Gamma(K_n,J_{\infty,\overline K/K}/\pi^i).
\end{align*}
最后一个等式来自
\[
 H^q(\Spec(\mathcal O_{\overline K}/\pi^i)/\Spec(\mathcal O_K/\pi^i)_{\mathrm{crys}},J^{[r]})=0
 \qquad(q\ge1,\ \text{任意 }r).
\]
参见 [Fo$_3$] §3 定理 1(i)，其中考虑的是 $K=K_0$ 的情形；同一证明也适用于一般情形。复合这些映射，得到 $\mu_{\xi,n}$。

将 $\mu_{r,\xi,n}$ 定义为复合
\begin{align*}
 S_n\otimes_{\mathcal O(G)}\Omega^1(G)
 &\xrightarrow{\mu_{\xi,n}}H^1(K_n,J_{\infty,\overline K/K}/J^{[r+1]}_{\infty,\overline K/K})\\
 &\xrightarrow{\cup\xi_n^{\otimes(-r)}}H^1\bigl(K_n,(T^{\otimes(-r)}\otimes_{\mathcal O_F}
 J_{\infty,\overline K/K}/J^{[r+1]}_{\infty,\overline K/K})/\pi^n\bigr).
\end{align*}
将 $\theta_{r,\xi,n}$ 定义为复合
\begin{align*}
 \Theta=\varprojlim_m\Omega^1(G)
 &\xrightarrow{(\mu_{r,\xi,m})_m}\varprojlim_m H^1\bigl(K_m,(T^{\otimes(-r)}\otimes_{\mathcal O_F}
 J_{\infty,\overline K/K}/J^{[r+1]}_{\infty,\overline K/K})/\pi^m\bigr)\\
 &\xrightarrow{\mathrm{trace}}\varprojlim_m H^1\bigl(K_n,(T^{\otimes(-r)}\otimes_{\mathcal O_F}
 J_{\infty,\overline K/K}/J^{[r+1]}_{\infty,\overline K/K})/\pi^m\bigr)\\
 &\xleftarrow{\sim}H^1\bigl(K_n,T^{\otimes(-r)}\otimes_{\mathcal O_F}
 J_{\infty,\overline K/K}/J^{[r+1]}_{\infty,\overline K/K}\bigr).
\end{align*}
\begin{statement}{定理 2.2.7}
对 $n,r\ge1$，复合映射
\begin{align*}
 \Theta&\xrightarrow{\theta_{r,\xi,n}}H^1\bigl(K_n,T^{\otimes(-r)}\otimes_{\mathcal O_F}
 J_{\infty,\overline K/K}/J^{[r+1]}_{\infty,\overline K/K}\bigr)\\
 &\longrightarrow H^1(K_n,V^{\otimes(-r)}\otimes_F B^1_{\mathrm{dR}}/B^{r+1}_{\mathrm{dR}})
 \simeq\operatorname{coLie}(G)^{\otimes r}\otimes_{\mathcal O_K}K_n
\end{align*}
将 $(\eta_m)_m\in\Theta$（各 $\eta_m\in\Omega^1(G)$）送到
\[
 -\frac{1}{(r-1)!}\pi^{-nr}\omega^{\otimes r}\otimes
 \left(\left(\frac d\omega\right)^{r-1}\left(\frac{\eta_n}{\omega}\right)\right)(\xi_n).
\]
\end{statement}
